% Folder 57, batch 09 (archivists' pages 161-180).
% Pass: Opus 5.5 (claude-opus-5-5), 2026-10-03 — first pass, unchecked against the pages by a human.
% Skipped: 162 (printed chapter opening, "Le langage des schémas"), 164 (typescript
% page "- 34 -", n° 297, on topological modules), 167 and 171 (blank).
\documentclass[11pt,a4paper]{article}
\input{../preamble/grothendieck}

\folder{57}
\batch{9}
\pages{161}{180}
\foldertitle{Picard : tapuscrit annoté (s.d.), lettres (s.d., 1962).}
\dating{1962-[vers 1968]}
\watermark{Édition de démonstration}

\begin{document}

\page{161}
\note{fragment ; le raisonnement commence avant ce feuillet}
\ill{} en correspondance \ill{} les $\mathcal{P}$ \uncertain{principaux} sur $X$ de groupe $GP(N)$, et les \struck{\ill{}} \ill{}

\[
\mathcal{P} \longrightarrow \mathcal{J}
\]

qui sont compatibles avec les opérations de $GP(N)$.
($GP(N)$ opérant sur $\mathcal{J}'$ \uncertain{dont les pts fixes} ($= \underline{\mathrm{Aut}}_{S}(X)$ \ill{}) \ill{}).

Donc $\simeq \Gamma(\mathcal{J}/GP(N))$ \uncertain{existe}, \ill{} \uncertain{correspond} \ill{}

\struck{\ill{}} \ill{} \uncertain{morphismes} $\mathcal{G} \to \mathcal{J}/GP(N)$ …

\marginal{\ill{} \ill{} $(X, P^{N})$ ; $\mathcal{J}' = \underline{\mathrm{Isom}}_{S}$ \ill{} et les \ill{} \emph{\uncertain{spéciales}}}

\page{163}
\note{notes de sa main, à l'encre, sur le verso d'une page dactylographiée sans rapport, qui n'est pas transcrite}
$Y \subset X$, \quad $X \times \mathrm{Pic}$

\begin{tikzcd}[column sep=small, row sep=small]
X \times \mathrm{Pic} \arrow[d] & \\
\mathrm{Pic} & \mathrm{Div} \arrow[l, "d"]
\end{tikzcd}

$\underline{L}^{Y}$

$\xi \in \underline{H}^{0}_{f}(\underline{L}^{Y})$ \ill{}

\[
\cdot \to \mathrm{Ker} \to \underline{H}^{0}_{f}(\underline{L}^{Y})^{\vee} \to \underline{\mathcal{O}}_{Y} \to 0
\]
\[
0 \to \underline{\mathcal{O}}_{S} \to \underline{H}^{0}_{f}(\underline{L}) \to \underline{H}^{0}_{f}(\underline{L}|Y)
\]
\[
\underline{J}\,\underline{L} \to \underline{L} \to \underline{L}|Y \to 0
\qquad\qquad \underline{L} = \underline{J}^{-1}
\]
\[
0 \to \underline{\mathcal{O}}_{X} \to \underline{L} \to \underline{L}|Y \to 0
\]
\[
0 \to \underline{\mathcal{O}}_{S} \to \underline{H}^{0}_{f}(\underline{L}) \to \underline{H}^{0}_{f}(\underline{L}|Y) \to \underline{H}^{1}_{f}(\underline{\mathcal{O}}_{X})
\]
sous les trois derniers termes : \uncertain{fibré projectif} \uncertain{de} $\mathrm{Div}/\mathrm{Pic}$ ; \uncertain{fibré tangent} $\mathrm{Div}/S$ ; \uncertain{tangent} $\mathrm{Pic}/S$.

\[
\to \underline{H}^{1}_{f}(\underline{L}) \longrightarrow \underline{H}^{1}_{f}(\underline{L}|Y) \to \underline{H}^{1}_{f}(\underline{\mathcal{O}}_{X})
\]
\note{sous les deux derniers termes, deux annotations dont on lit seulement des mots : « … obstruction \uncertain{universelle} identique \uncertain{pour} $\mathrm{Div}/\mathrm{Pic}$ … et \ill{} \ill{} \ill{} » ; « … \uncertain{lisse} … $\mathrm{Div}$ … »}

\page{165}
\section{AK — Pseudomorphismes}
\note{« AK » en tête à gauche, « Pseudomorphismes » souligné en tête à droite, au crayon ; feuillet en largeur}

$X$ un préschéma.
\note{ci-dessus et dans la suite, les retours à la ligne sont ceux du feuillet}
Ouvert schématiquement \struck{\ill{}} dense $U$ : si $X$ est le plus \struck{petit} \uncertain{sous}-préschéma fermé \ill{} $X$ dominant $U$, i.e. si $\overline{U} = X$ [lorsque $U \to X$ est quasi-compact \ill{} \ill{} \ill{} $V$, $U \cap V$ est schématiquement dense dans $X$ --- aussi \ill{} \ill{} \ill{} mieux \ill{} \ill{} définition : toute \struck{\uncertain{fonction}} \add{\uncertain{section} de $\underline{\mathcal{O}}_{X}$} sur un \ill{} de $X$ \ill{} sur $U \cap V$ est nulle].

\struck{\ill{}} Condition locale \ill{} \ill{} \uncertain{section} [si $X = \bigcup X_{i}$, les $X_{i}$ \ill{}, alors \ill{} les $U \cap X_{i}$ sch. dense dans les $X_{i}$].

$U$ schématiquement dense dans $X \Longleftrightarrow$ \ill{} \ill{} $V$ q.c. $U \subset V \subset X$.
$U$ schém. dense dans $X \Rightarrow U$ schém. dense dans \ill{} \ill{}.
Si $U, U'$ sont schém. denses, alors $U \cap U'$ l'est.

\marginal{\ill{} \ill{}}

[Soit $X$ et $Y$ sur $S$, \struck{$X$} \add{$Y$} \uncertain{séparé} sur $S$, \struck{$X \to$} $U$ schém. dense dans $X$.
Si $f, g : X \to Y$ sont \ill{} $S$-morphismes tels que $f|U = g|U$ \ill{}]

\struck{\ill{}} \add{\uncertain{Pseudo}}-morphismes : classe d'équivalence \struck{d'\uncertain{applications}} de morphismes définis sur des \ill{} schématiquement denses, \uncertain{équivalents} \ill{} $S$-morphismes. Pour un \ill{} $S$-morphisme dans un $S$-\uncertain{préschéma séparé}, il y a un \add{(plus grand ouvert)} domaine de définition.
\struck{\ill{} $D(f|V) =$} $\mathrm{Dom}(f|V) = \mathrm{Dom}(f) \cap V$

\page{166}
\textbf{Ex.} Si $X$ est \add{réduit, alors \uncertain{schématiquement} dense $=$ dense} \struck{\ill{} \ill{} \ill{} \ill{} \ill{} \ill{} irréductibles}
pseudo-morphismes $=$ applications rationnelles.

Pour tout ouvert $U$ de $X$, soit $F(U)$ l'ens. des pseudo-$S$-morphismes de $U$ dans $Y$ ($S$-schéma fixé). Applications de restriction $F(U) \to F(V)$ si $U \supset V$. Cela \struck{est} \uncertain{définit} un \emph{faisceau} [grâce à la relation $\mathrm{Dom}\, f_{i}|U_{ij} = \mathrm{Dom}(f_{i}) \cap U_{ij}$] \uncertain{appelé} \emph{faisceau des} \struck{\ill{}} \emph{pseudo-$S$-morphismes} de $X$ dans $Y$.

En particulier, on obtient ainsi le faisceau des pseudo-sections de $\mathcal{O}_{X}$ (faisant $X = S$, $Y = X[t]$) qui est un faisceau d'anneaux \struck{\ill{}} \uncertain{contenant} $\underline{\mathcal{O}}_{X}$, en particulier un faisceau d'algèbres.

\struck{Prop. \ill{} \ill{} faisceau est quasi-cohérent \ill{} \ill{} \ill{} \ill{} $X = \mathrm{Spec}(A)$.}
\note{les deux dernières lignes sont biffées}

\page{168}
\marginal{Diviseurs de fonctions}
Soit $X$ une variété irréductible \emph{normale} déf/$k$
$K = k(X)$ le corps des fonctions rationnelles/$k$ sur $X$.
Soit $V$ un \struck{\ill{} \ill{} irréductible} \add{diviseur premier déf.} sur $k$ sur $X$.

Alors l'ensemble des $f \in K$ \emph{définies} en $V$ forment un anneau de valuation discrète, \struck{\ill{}} $\mathcal{O}_{(V)}$.

En effet, c'est un \struck{\ill{}} localité normale, dont la dimension de Krull est 1.

Soit $v_{V}$ la valuation \add{de $K$} \uncertain{attachée} à cet \ill{} \ill{} \ill{} valuation de $K$. Appelons \emph{diviseur} \add{déf/$k$} de $X$ toute \ill{} \ill{} formelle des diviseurs premiers \uncertain{rat.}\ déf/$k$. Posons, pour tout $f \in K$
\[
\mathrm{div} f = \sum_{\substack{V \text{ div.} \\ \text{pr}/k}} v_{V}(f) \cdot V
\]
Cette somme a un sens, car pour $f$ donnée, il n'y a qu'un nb fini de termes non nuls ; \struck{\ill{}} pour le voir, on se ramène au cas \struck{\ill{}} affine, \uncertain{alors} les diviseurs pr. déf/$k$ correspondent aux idéaux premiers \ill{} \ill{} minimaux de l'\ill{} \ill{} \uncertain{normal} noethérien $A = k[X]$, et le \uncertain{résultat} est bien connu.

\textbf{Théorème} \struck{\ill{} \ill{} \ill{} équivalent} $\mathrm{div} f \geqslant 0 \Longleftrightarrow f$ est partout définie

C'est \uncertain{résultat} de la théorie des \ill{} \ill{} \uncertain{normaux} :
$A = \bigcap A_{\mathfrak{p}}$ ($\mathfrak{p}$ premier minimal) si $A$ est normal.
En effet, \ill{} \ill{} \ill{} \ill{} \ill{} \ill{} \uncertain{définition},
et $\mathrm{div} f \geqslant 0 \Longleftrightarrow f \in \bigcap A_{\mathfrak{p}}$.

\page{169}
\textbf{Corollaire 1} \struck{Pour que $f$ soit défini en un \ill{} \ill{} sous-variété irréductible $Y$, il faut et il suffit que} \add{L'ens. \uncertain{des points} de $X$ où $f$ \ill{} \ill{} défini se} \struck{\ill{} le voisinage des \ill{} \ill{} $V$ tels que} $v_{V}(f) < 0$. \struck{\ill{} \ill{} \ill{} $V$ \ill{}} \ill{} \ill{} \ill{} si $f$ est une application rationnelle de $X$ dans une variété \struck{\ill{}} ($X$ normale) \uncertain{c'est que} \ill{} une sous-variété pure de dim $n-1$.

\struck{\ill{}} N.B. Les résultats précédents \uncertain{restent} \ill{} \ill{} \uncertain{essentiel} si \struck{\ill{}} \add{l'ens. des pts singuliers de $X$} est de dim $\leqslant n-2$.

\textbf{Th.} Soit $X$ une variété \struck{\ill{} \ill{}} \add{normale de dim $n$}, $Y$ une variété complète, $f : X \to Y$ une application rationnelle. Alors l'ens. des pts où $f$ \ill{} \uncertain{pas} définie est de dim $\leqslant n-2$.

\struck{\ill{} \ill{} \ill{} $X$}
\uncertain{Deux démonstrations} :

a) \ill{} \ill{} \ill{}, en regardant le graphe de $f$, que $f$ est \uncertain{inverse} d'un \emph{morphisme} birationnel $g : Y \to X$, et $g$ est \struck{\ill{}} propre.
\struck{\ill{}} Soit $X'$ l'un des pts \struck{\ill{}} \ill{} tels que $g^{-1}(x)$ soit de dim $> 0$, alors $X'$ est de dim $\leqslant n-2$, car sinon $g^{-1}(X')$ serait de dim $\geqslant (n-1)+1 = n$, absurde. \struck{Donc $g^{-1}(X)$} D'\ill{} \uncertain{part}, il résulte du Main th. que \ill{} dehors de $\complement X'$, $g$ est un isom. biunivoque, donc $g^{-1}$ est défini sur $\complement X'$. O.K.

N.B. \uncertain{Bien entendu}, le th. se généralise \ill{} \ill{} au cas où \struck{\ill{}} l'ens. des \uncertain{pts singuliers} de $X$ est de dim $\leqslant n-2$

\page{170}
\note{le raisonnement commence avant ce feuillet}
\ill{} \ill{} $X$ irréductible, $f$ \uncertain{domine} \ill{}, \ill{} \uncertain{précisément} \ill{}.

b) Soit $V$ un hyperplan \uncertain{irréductible} de $X$. \ill{} voir que $f$ est définie en \uncertain{un} \ill{} au \uncertain{moins} \ill{} de $V$, \struck{\ill{}} i.e. $f$ « définie en \ill{} localité $\mathcal{O}_{(V)}$ », il faut montrer qu'il existe une localité de $Y$ dominée par $\mathcal{O}_{(V)}$. \uncertain{Or} \ill{} résulte de la définition valuative des variétés complètes, et du fait que $\mathcal{O}_{(V)}$ est un anneau de valuation.

\page{172}
\section{Le foncteur de Picard}
\marginal{Faits élémentaires sur les classes de diviseurs, faisceaux inversibles, immersions projectives etc.}

Soit $X$ un $S$-préschéma, on pose

$\text{Prépic}(X/S) =$ \struck{$H^{1}(X, \underline{\mathcal{O}}_{X}^{*})$} $\mathrm{Pic}(X)$ \struck{\ill{}} $= H^{1}(X, \underline{\mathcal{O}}_{X}^{*})$
\[
\underline{\underline{\text{Prépic}}}_{X/S}(S') = \text{Prépic}(X'/S') = H^{1}(X', \underline{\mathcal{O}}_{X'}^{*})
\qquad \text{où } X' = X \times_{S} S'
\]
donc $\underline{\underline{\text{Prépic}}}_{X/S}$ est bien un foncteur \emph{contravariant} en le préschéma $S'$.
\[
\text{Prépic}'(X/S) = \Gamma(S, \underline{H}^{1}_{X/S}(\underline{\mathcal{O}}_{X}^{*}))
\]
\[
\underline{\underline{\text{Prépic}}}'_{X/S}(S') = \Gamma(S', \underline{H}^{1}_{X'/S'}(\underline{\mathcal{O}}_{X'}^{*}))
\qquad \text{où } X' = X \times_{S} S'
\]
\uncertain{On passe} de $\underline{\underline{\text{Prépic}}}$ à $\underline{\underline{\text{Prépic}}}'$ en \uncertain{remplaçant} la \ill{} \uncertain{grossière} \ill{} \ill{} \uncertain{recouvrements} \ill{}.

\textbf{Proposition 1} \uncertain{Supposons} $\underline{H}^{0}_{X/S}(\underline{\mathcal{O}}_{X}) = \underline{\mathcal{O}}_{S}$. Alors \ill{} un \struck{\uncertain{isomorphisme}} \add{\uncertain{monomorphisme}} canonique
\[
\text{Prépic}'(X/S) \leftarrow \text{Prépic}(X/S)/\mathrm{Pic\,abs}(S)
\]
qui est un isomorphisme si $X/S$ \uncertain{admet} une section.
\note{la flèche, tracée en gras, va de droite à gauche}
Résulte de l'\uncertain{interprétation} \uncertain{plus} \uncertain{précise} \ill{}

\textbf{Corollaire 1} Le foncteur $\underline{F} \to \underline{F} \otimes_{S} X$ des faisceaux loc. libres sur $S$ dans les faisceaux loc. libres sur $X$ \struck{\uncertain{définis}} \add{est} pleinement \emph{\uncertain{fidèle}}, \uncertain{i.e.} \ill{} \ill{} \ill{} est \uncertain{fermée} \ill{} \ill{} faisceaux loc. libres sur $X$ qui sont \uncertain{principaux} loc. \ill{} \ill{} \ill{} de $Y$, \ill{} foncteur inverse \ill{} \ill{} \ill{} est le foncteur $\underline{H}^{0}_{X/S}$.

\marginal{\uncertain{expliciter} \ill{} \ill{} \ill{} \ill{} \ill{} $\mathrm{Pic}(X/S) \to \text{Prépic}'(X/S)$ ; \ill{} \struck{\uncertain{corollaire}} \ill{} \ill{} d'une \uncertain{suite} exacte \uncertain{évidente} ; \ill{} \ill{} \ill{} \ill{} \uncertain{préliminaires}}

\page{173}
Il en résulte en \uncertain{plus} que l'on a une \uncertain{suite} exacte (\ill{} \ill{} \ill{} \ill{} \ill{} \ill{} faisceaux \uncertain{inversibles})
\[
0 \to \mathrm{Pic}(S) \to \mathrm{Pic}(X) \to \text{Prépic}'(X/S)
\]
\struck{l'image} \add{le \uncertain{dernier}} \uncertain{étant} \uncertain{formé} des sections de $\underline{H}^{1}_{X/S}(\underline{\mathcal{O}}_{X}^{*})$ qui \ill{} \ill{}, \uncertain{l'image} de $\underline{H}^{1}(\underline{\mathcal{O}}_{X}^{*})$ \uncertain{i.e.}\ dans l'\uncertain{obstruction} dans $H^{2}(S, \underline{\mathcal{O}}_{S}^{*})$ est nulle,

[N.B. \ill{} \ill{} \ill{} \ill{} \ill{} \ill{} suite exacte \ill{} \ill{} \ill{} \ill{} \ill{} \ill{} \ill{} \ill{} \uncertain{termes} de \ill{} \ill{} \ill{} de Leray
\[
0 \to H^{1}(S, \underline{H}^{0}(\underline{\mathcal{O}}_{X}^{*})) \to H^{1}(X, \underline{\mathcal{O}}_{X}^{*}) \to H^{0}(S, \underline{H}^{1}(\underline{\mathcal{O}}_{X}^{*})) \to H^{2}(S, \underline{H}^{0}(\underline{\mathcal{O}}_{X}^{*}))
\]
\[
H^{2}(S, \underline{H}^{0}(\underline{\mathcal{O}}_{X}^{*})) \to H^{2}(X, \underline{\mathcal{O}}_{X}^{*})
\]
\note{la dernière flèche descend verticalement sous $H^{2}(S, \underline{H}^{0}(\underline{\mathcal{O}}_{X}^{*}))$}
\uncertain{compte tenu} \ill{} \ill{} $\underline{H}^{0}(\underline{\mathcal{O}}_{X}^{*}) = \underline{\mathcal{O}}_{S}^{*}$]. S'il \uncertain{existe} une section de $X/S$, alors \ill{} \struck{\ill{}}
\struck{\ill{} \ill{} \ill{} \ill{}}
\ill{} \uncertain{résultat} \uncertain{plus complet} \ill{}

\marginal{[\uncertain{Interprétation} si $X/S$ a une section] ; \ill{} \ill{} \ill{} \ill{} \ill{} \ill{} \ill{} \ill{} l'existence \ill{} \ill{} $\underline{H}^{0}_{X/S}(\underline{\mathcal{O}}_{X}) = \underline{\mathcal{O}}_{S}$}

\textbf{Corollaire 2} Soit $\underline{C}$ la catégorie des faisceaux \ill{} sur $X$ munis d'un \emph{\uncertain{isomorphisme}} $\underline{\mathcal{O}}_{S} \to s^{*}(\underline{L})$. Alors (i) \uncertain{tout} \ill{} \uncertain{automorphisme} \ill{} \ill{} \ill{} de $\underline{C}$ est trivial (ii) \uncertain{L'ensemble} \struck{\ill{}} \ill{} de $\text{Prépic}(X/S)$ \uncertain{provient} d'un $\underline{L} \in \underline{C}$, déterminé à isom. \uncertain{unique} près.

\page{174}
Pour prouver (ii), on considère un recouvrement de $S$ par des $U_{i}$ tels que $\xi$ soit donné par des $\underline{L}_{i}$ sur $X_{i} = X|U_{i}$, prenant les $U_{i}$ assez petits

\struck{[\ill{} \ill{} \ill{} $S$…] \ill{} \ill{} \ill{} \ill{}}

Remplaçons $\underline{L}_{i}$ par $\underline{L}_{i} \otimes_{S} s^{*}(\underline{L}_{i})^{-1}$, on \ill{} \uncertain{supposer} \ill{} \ill{} \ill{} \uncertain{isomorphismes} $\mathcal{O}_{S}|U_{i} \xrightarrow{\varphi_{i}} s^{*}(\underline{L}_{i})$. Dans $X_{ij} = X|U_{i} \cap U_{j}$, $\underline{L}_{i}$ et $\underline{L}_{j}$ sont \emph{$U_{ij}$-équivalents}, donc il y a un isomorphisme \struck{\ill{}} \ill{} \ill{} de $\underline{L}_{i}$ sur $\underline{L}_{j}$ compatible avec les $\varphi_{j}$ et $\varphi_{i}$. \struck{\ill{}} L'\uncertain{assertion} résulte que les \ill{} \uncertain{satisfont} la condition de recollement, d'où le \uncertain{résultat}. Question : donner une formulation \uncertain{cohomologique} du raisonnement.

\marginal{\uncertain{démontrer} \ill{} \ill{} \ill{} il \ill{} \ill{} \ill{} \ill{} \ill{} \ill{} \ill{} (\ill{} \ill{})}

\textbf{\add{Proposition 2}} \struck{\ill{}} \ill{} \ill{} \ill{} \ill{} $\underline{H}^{0}_{X/S}(\underline{\mathcal{O}}_{X}) = \underline{\mathcal{O}}_{S}$, et $X$ \uncertain{quasi-compact} \uncertain{séparé} sur $S$, \ill{} \ill{} \ill{} \ill{} \ill{} \ill{} \ill{} \ill{} \uncertain{extension} \uncertain{plate} de la base. Soit $S' \to S$ une extension \add{\uncertain{quasi-compacte}} fid. plate de la base. (i) Alors
\[
\text{\struck{$\mathrm{Pic}$}}\ \text{Prépic}'(X/S) \to \text{Prépic}'(X'/S')
\]
est injectif. (ii) Si $X/S$ \uncertain{localement} \struck{\ill{}} \ill{} une section, \ill{} \ill{} \ill{} \ill{}

\page{175}
\[
\to \text{Prépic}'(X/S) \to \text{Prépic}'(X'/S') \rightrightarrows \text{Prépic}'(X''/S'')
\]
[où $S'' = S' \times_{S} S'$] est exacte \struck{\ill{} \ill{} \ill{}}

\struck{$\underline{\underline{\text{Prépic}}}_{X/S}$}

\struck{\ill{}} En effet, le point (ii) est \ill{} \ill{} conséquence immédiate de

\textbf{Corollaire 1} sous les conditions de prop. 1, \marginal{$X$ \uncertain{séparé} \uncertain{quasi-compact} sur $S$}

Corollaire 1, soit $S' \to S$ une extension \add{\uncertain{quasi-compacte}} fid. plate de la base. Soit $\underline{E}$ \struck{\ill{}} \add{loc. libre} \struck{\ill{}} sur $X$. Pour qu'il soit isomorphe \struck{\ill{}} à l'image inverse d'un faisceau loc. libre sur $S$, il faut que $\underline{E}'$ soit isomorphe à l'image inverse d'un faisceau loc. libre sur $S$.

En effet, il faut exprimer que $f_{*}(\underline{E})$ est loc. libre et \struck{\ill{}} $f^{*}f_{*}(\underline{E}) \to \underline{E}$ \ill{} isom. Or $f_{*}(\underline{E})' = f_{*}(\underline{E}')$ \add{(\uncertain{car} \uncertain{fid.} \ill{} \ill{})} \ill{} \ill{} \uncertain{ramène} à des \uncertain{assertions} \uncertain{évidentes} de stabilité de propriétés \uncertain{par} \ill{} fid. plates (\uncertain{Chap.}\ IV).

\page{176}
Pour prouver (ii), on \ill{} \uncertain{ramené} à prouver \ill{} \ill{} \ill{} de $\text{Prépic}'(X'/S')$ dont les \uncertain{images} sur les \ill{}, provient d'un \ill{} de $\text{Prépic}'(X/S)$. À cause de l'injectivité, on voit \uncertain{que} \ill{} \ill{} sur $S$, on \ill{} \uncertain{donc} supposer que \ill{} \ill{} une section de $X$ sur $S$. \ill{} \ill{} \ill{} identifier les \struck{\uncertain{objets}} \add{\uncertain{éléments}} de $\text{Prépic}(X/S)$ \add{(\uncertain{resp.}\ \ill{} \ill{})} \ill{} des faisceaux inversibles sur $X$ \uncertain{à} section \uncertain{rigidifiée} sur $S$, \uncertain{on} \uncertain{peut} \ill{} \uncertain{appliquer} la théorie de la descente \emph{fid. plate}.

\textbf{Corollaire 2} \uncertain{Lorsque} \uncertain{combine} \ill{} Prop. (ii) \ill{} \ill{} \ill{} \add{\struck{\ill{}} \uncertain{faisceaux} sur $X$}.

\textbf{Corollaire 3} Supposons que la relation $\underline{H}^{0}_{X/S}(\underline{\mathcal{O}}_{X}) = \underline{\mathcal{O}}_{S}$ \uncertain{reste} vraie \uncertain{par} \uncertain{toute} extension de la base (i.e., $X$ \uncertain{est} plat sur $S$ \ill{} \ill{}, \add{et $X/S$ a localement une section}, et $\underline{H}^{0}_{X/S}(\underline{\mathcal{O}}_{X}) = \underline{\mathcal{O}}_{S}$), (Alors le foncteur $\underline{\underline{\text{Prépic}}}'_{X/S}$ est \uncertain{compatible} [$=$ exact \uncertain{pour}] \ill{} la descente fidèlement plate \uncertain{quasi-compacte}.

\page{177}
\struck{\ill{}}

$\underline{\underline{\mathrm{Pic}}}_{X/S} =$ enveloppe exacte de $\underline{\underline{\text{Prépic}}}_{X/S}$ (ou $\underline{\underline{\text{Prépic}}}'_{X/S}$) pour la descente fid. plate quasi-compacte.

[Il y a un \uncertain{ennui}, car un \emph{\uncertain{recouvrement}} \ill{} \ill{} d'un préschéma \ill{} \uncertain{quasi-compact} \ill{} \ill{} \ill{} \ill{} \uncertain{majoré} par un morphisme fid. plat qu.-compact. Il \uncertain{faut} \uncertain{donc} \uncertain{dire} \ill{} \ill{} \ill{} fid. plat qu.-compact : \struck{fid. plat qu.} \uncertain{après} \struck{\uncertain{localisation}} \emph{morphisme de localisation fidèlement plate}]

Donc on a des morphismes fonctoriels
\[
\underline{\underline{\text{Prépic}}}_{X/S} \xrightarrow{\alpha} \underline{\underline{\text{Prépic}}}'_{X/S} \xrightarrow{\beta} \underline{\underline{\mathrm{Pic}}}_{X/S}
\]
\note{au-dessus du premier « Pic » de $\text{Prépic}$, une surcharge illisible}
Le premier morphisme n'est pas injectif en général [\uncertain{noyau} \ill{}, si $\underline{H}^{0}_{X/S}(\underline{\mathcal{O}}_{X}) = \underline{\mathcal{O}}_{S}$, et \ill{} \ill{} reste vrai par changement de base] est le foncteur $\underline{\mathrm{Pic}}$ (absolu) \add{(\ill{} \ill{} \ill{} \ill{} \ill{}…)} en \uncertain{vertu} de prop. 1, (i), [Il est surjectif si \ill{} \uncertain{suppose} en plus de la condition \ill{} [ ] que $X/S$ a une section. \struck{\ill{}}

\page{178}
D'autre part $\beta$ est un \uncertain{monomorphisme}, \ill{} \uncertain{l'hypothèse} $\underline{H}^{0}_{X/S}(\underline{\mathcal{O}}_{X}) = \underline{\mathcal{O}}_{S}$ « universelle » ; et il est surjectif dès qu'il \uncertain{l'est} si on suppose \emph{de plus} que \struck{\ill{}} $X$ a localement une section sur $S$.

\textbf{\uncertain{Exemple}} où $\beta$ n'est pas surjectif : variété de Severi-Brauer non triviale \ill{} \ill{} \ill{} \ill{} \ill{} \ill{} \ill{} \ill{} \ill{} \ill{} \ill{} \ill{} \ill{} \uncertain{distribution} : la surjectivité de $\beta$ \ill{} \ill{} \ill{} \ill{} \ill{} $H^{2}(S, \mathbf{G}_{m})$…

\textbf{Remarques} On voit donc que \struck{le} foncteur $\underline{\underline{\text{Prépic}}}$ n'est \add{\ill{} \ill{} \ill{} \ill{}} pas représentable \struck{en général, \ill{} \ill{}} \struck{par} \ill{} \uncertain{triviaux} \ill{} \ill{} ; $\underline{\underline{\text{Prépic}}}'_{X/S}$ ne \ill{} \ill{} être représentable que si $\beta$ est un isom. et si $\underline{\underline{\mathrm{Pic}}}_{X/S}$ est représentable, \uncertain{ce qui} \struck{\ill{}} [\ill{} \ill{} \ill{} \ill{} \uncertain{pratique} \ill{} \ill{} \uncertain{des} conditions, \ill{} \ill{} \ill{} \ill{}, \ill{} \ill{} l'existence d'une section de $X$ sur $S$]. \uncertain{Question} : $\underline{\underline{\mathrm{Pic}}}_{X/S}$, il est représentable dans des cas \uncertain{suffisamment} importants \uncertain{qui} seront \uncertain{étudiés} \uncertain{par} la suite. C'est \ill{} \ill{} \ill{} le \uncertain{plus} important \uncertain{pour} \ill{} \uncertain{propos}.

\page{179}
$X$ \struck{\uncertain{plat}} \ill{} \ill{} de $S$

$Y$ sous-schéma divisoriel de $X$ transversal aux fibres est défini par un faisceau \emph{inversible} $\underline{L}$ sur $X$ et une section $s$ de $\underline{L}$ \struck{\ill{}} non diviseur de 0 sur les fibres.

Partons d'un faisceau inversible $\underline{L}$ sur $X$, et déterminons les \struck{\uncertain{faisceaux}} sous-schémas divisoriels de $X$ transversaux aux fibres qui définissent \ill{} $\underline{L}(Y)$ \struck{\ill{}} qui sont \struck{\ill{} \ill{}}
\struck{\ill{} \ill{}} isomorphes [resp. $S$-équivalents, resp.
\struck{$Y$-équivalents}] \struck{\ill{}} \uncertain{Pic}-équivalents \uncertain{rel.}\ à $Y$]
à $\underline{L}$. // Les premiers sont ceux qui se définissent par les sections de $\underline{L}$ non diviseurs de 0 sur les fibres, sections qu'il faut prendre \uncertain{modulo}
$\Gamma(S, f_{*}(\underline{\mathcal{O}}_{X}^{*}))$ ; si $f_{*}(\underline{\mathcal{O}}_{X}) = \underline{\mathcal{O}}_{S}$, \ill{} les \ill{} donc modulo $\underline{\mathcal{O}}_{S}^{*}$, donc on \ill{} la \uncertain{partie} \struck{de $\underline{H}^{0}_{f}(\underline{L})$} \add{\uncertain{convenable}} d'une \ill{}

$\underline{H}^{0}_{f}(\underline{L})^{*}$ de $\underline{H}^{0}_{f}(\underline{L})$ \uncertain{par} le \uncertain{groupe}

$\underline{H}^{0}_{f}(\underline{\mathcal{O}}_{X}^{*}) = \underline{H}^{0}_{f}(\underline{\mathcal{O}}_{X})^{*}$, donc \uncertain{dans} le \uncertain{cas} \uncertain{favorable} \uncertain{par} $\underline{\mathcal{O}}_{S}^{*}$. --- Dans le deuxième cas, on trouve les \struck{\ill{}} \add{sections} de $\underline{H}^{0}_{f}(\underline{L})^{*}/\underline{H}^{0}_{f}(\underline{\mathcal{O}}_{X})^{*}$ dans le \uncertain{cas} le \uncertain{favorable} de $\underline{H}^{0}_{f}(\underline{L})^{*}/\underline{\mathcal{O}}_{S}^{*}$. ---

\page{180}
Il en est encore de \struck{\ill{}} \struck{\ill{}} \ill{} la troisième, si on suppose que $X$ est quasi-compact séparé sur $S$ et que $\underline{H}^{0}_{f}(\underline{\mathcal{O}}_{X}) = \underline{\mathcal{O}}_{S}$ --- \ill{} \ill{} \uncertain{pas} la $S$-équivalence \uncertain{est} \uncertain{identique} : la Picard-équivalence \uncertain{rel.}\ à $Y$

Supposons \uncertain{maintenant} que
$X$ quasi-compact séparé sur $S$, $\underline{H}^{0}_{f}(\underline{\mathcal{O}}_{X}) = \underline{\mathcal{O}}_{S}$ et Kif-kif \uncertain{par} \add{\uncertain{extension} de la base}

et que $\underline{\underline{\mathrm{Pic}}}_{X/S}$ et $\underline{\underline{\mathrm{Div}}}_{X/S}$ existent. Alors l'ensemble précédent
\[
H^{0}(S, \underline{H}^{0}_{f}(\underline{L})^{*}/\underline{\mathcal{O}}_{S}^{*})
\]
peut s'interpréter comme l'ensemble des sections $\theta$ de $\underline{\underline{\mathrm{Div}}}$ sur $S$ telles que $p\theta = \lambda$, où $\lambda$ est la section de $\underline{\underline{\mathrm{Pic}}}$ sur $S$ qui correspond à $\underline{L}$. Donc, posons $\underline{\underline{\mathrm{Div}}}^{(L)} = \varphi^{-1}(\lambda(S))$ [sous-préschéma de $\underline{\underline{\mathrm{Div}}}$] c'est l'ens. des sections de $\underline{\underline{\mathrm{Div}}}^{L}_{X/S}$ sur $S$. \struck{\ill{}} Ceci est compatible \ill{} extension de la base $S'$, donc le foncteur
\[
\underline{\underline{\mathrm{Div}}}^{L}_{X/S}(S') = H^{0}(S', \underline{H}^{0}_{f'}(\underline{L}')^{*}/\underline{\mathcal{O}}_{S'}^{*})
\]
\ill{} $S'$ est représentable. \struck{\uncertain{Ceci dit}} [La \uncertain{définition} \ill{} \ill{} condition nécessaire intéressante \ill{} \uncertain{termes} des diviseurs, pour l'existence

\note{la phrase se poursuit au-delà de ce feuillet}

\begin{tikzcd}[column sep=small, row sep=small]
\underline{\underline{\mathrm{Pic}}} \arrow[dr, "\varphi"'] & & \underline{\underline{\mathrm{Div}}} \arrow[ll, "p"'] \arrow[dl, "\psi"] \\
& S &
\end{tikzcd}

\note{ce petit diagramme est tracé dans la marge gauche, à hauteur de « $p\theta = \lambda$ »}

\end{document}
